problem:  
Let \((M_i^n, g_i)\) be a sequence of complete, noncompact Riemannian manifolds with **nonnegative sectional curvature** and \(|\mathrm{sec}_{g_i}|\le1\), converging in the pointed Gromov–Hausdorff sense to a lower-dimensional limit space \((X,p_\infty)\).  
For each \(M_i\), let \(\Sigma_i\subset M_i\) denote the soul, so that \(\pi_1(M_i)\cong\pi_1(\Sigma_i)\) is **virtually abelian**.  

Suppose further that:  
1. The souls \(\Sigma_i\) have a uniform dimension bound and uniformly bounded normal bundle curvature;  
2. The diameter of \(\Sigma_i\) is uniformly bounded above.  

**Question:**  
Under these geometric constraints, must the abelian rank \(r_i = \operatorname{rank}(\pi_1(M_i))\) stabilize for large \(i\)?  
Equivalently, can a collapsing sequence of nonnegatively curved manifolds with bounded sectional curvature and *uniformly controlled souls* exhibit infinitely oscillating abelian rank?  

Why is it a "good" problem:  
This version introduces precise, natural geometric control (bounded soul geometry and normal bundle curvature) that is not covered by existing collapse theory. The interaction between the **soul structure**, which encodes the stable topological core, and **bounded-curvature collapse**, governed by the Cheeger–Fukaya–Gromov theory, is largely unexplored under such uniformity hypotheses.  

If rank stabilization holds, it would identify a new stability phenomenon linking the topology of souls to the ambient geometry — a subtle rigidity principle for nonnegative curvature. If not, it would reveal hidden flexibility in collapse compatible with these additional geometric bounds.  

The problem is succinct, uses only elementary notions (rank of a group, sectional curvature, soul, Gromov–Hausdorff convergence), yet points directly toward a frontier joining Riemannian geometry, geometric topology, and group theory — a hallmark of a “good” research problem.